\section{What the test does not give}\label{sec:limits}

The entries of $S$ count walks, not paths, so only their positivity is used: on $G_1$, $S_{22} = 2$ counts the walks $n_2, n_1, n_2$ and $n_2, n_3, n_2$, but only one path, $(n_2)$, joins $n_2$ to itself (\Cref{ex:board}, \Cref{ex:test1}). The test is also not the cheapest way to decide connectivity. It forms $N - 2$ matrix products, while a graph search from one vertex such as depth-first search runs in time $O(E)$ for $E$ edges, since it examines each edge at most twice \cite{wayne_cos226_undirected}. That the test is therefore the more expensive one is an interpretation: the source states the running time for depth-first search only, says nothing about breadth-first search or about this test, and no complexity analysis is made here.

Two sources state these results with the word ``paths'' where walks are meant. Paramadevan and Sotheeswaran's Theorem 1 says the entry of $A^l$ is the number of ``different paths'' of length $l$ \cite[Theorem 1, p.~16]{paramadevan2021adjacency} and the paper never defines ``path''; read as walks it is \Cref{thm:walks}, read with the paths of \Cref{def:path} it fails, since $(A^2)_{11} = 1$ on $G_1$ while no path of length $2$ joins $n_1$ to itself. Bai's proof sketch says $(I+A)^{n-1}$ ``aggregates all paths of length $\le n-1$'' \cite[p.~2]{bai_bimsa_adjacency}, while his Theorem 1 says that powers of $A$ count walks, not paths \cite[p.~1]{bai_bimsa_adjacency}; the entries of $(I+A)^{n-1}$ count walks too (\Cref{rem:binom}).
